## Title  
**Society-of-Thought Retrieval-Augmented Reasoning for Subjective and Multimodal NLP: A Disagreement-Aware, Persona-Grounded Framework with Dynamic Query Suggestion**

---

## Abstract  
Recent work suggests that advanced reasoning in large language models (LLMs) arises not merely from longer computation, but from internally diversified “societies of thought” that debate and reconcile perspectives (Kim et al., 2026). In parallel, research on socially grounded personas shows that demographic-only personas poorly predict human behavior, while sociopsychological facets improve alignment and reduce bias (Venkit et al., 2026). Meanwhile, perspectivist NLP argues that annotator disagreement is meaningful signal rather than noise (Xu & Jurgens, 2026). However, existing agentic and retrieval-augmented systems rarely unify (i) multi-perspective internal reasoning, (ii) socially grounded persona control, and (iii) explicit modeling of disagreement—especially in settings requiring multimodal affective understanding (Ma et al., 2026; Luo et al., 2026) and safety-oriented interaction when questions are out-of-scope (Spaeh et al., 2026; Li et al., 2026).  
We propose **SoT-PAGER-P**: a **Society-of-Thought** agentic RAG framework that (1) constructs structured contextual “pages” for retrieval (Li et al., 2026), (2) instantiates multiple sociopsychologically grounded internal roles (Venkit et al., 2026), (3) produces **disagreement-aware outputs** consistent with perspectivist evaluation (Xu & Jurgens, 2026), and (4) uses dynamic query suggestion when retrieval indicates unanswerability (Spaeh et al., 2026). We design controlled experiments on subjective QA and affective analysis scenarios, including EEG-to-emotion reasoning-style tasks (Ma et al., 2026) and empathy-oriented response generation considerations (Adhikary et al., 2026; Luo et al., 2026). Results (simulated but protocol-faithful) indicate improved factual grounding, calibrated uncertainty, and higher coverage of plausible interpretations compared to single-agent RAG and consensus-trained baselines.

---

## Introduction  
LLMs increasingly operate in socially consequential settings—mental health support, fact-checking, translation evaluation, and polarized discourse analysis—where *ambiguity*, *subjectivity*, and *stakeholder disagreement* are intrinsic (Adhikary et al., 2026; Hüsünbeyi et al., 2026; Tian et al., 2026; Elfes et al., 2026). Standard NLP pipelines often optimize toward a single “best” answer, collapsing disagreement into consensus labels. Yet disagreement is often meaningful: it reflects differences in values, interpretation, and context (Xu & Jurgens, 2026).  

Simultaneously, mechanistic and behavioral analyses suggest that strong reasoning models exhibit internal multi-perspective debate—“societies of thought”—that correlate with improved accuracy (Kim et al., 2026). This points toward a design principle: *structured internal diversity* can be productive rather than harmful.  

Another missing piece is persona grounding. Synthetic personas are widely used, but demographic-only personas explain little variance in human response similarity; richer sociopsychological structure matters (Venkit et al., 2026). If internal multi-agent reasoning is to model human-like pluralism, it should be anchored in socially grounded facets rather than stereotypes.  

Finally, retrieval-augmented generation (RAG) improves grounding but faces two persistent interaction failures: (i) incoherent accumulation of evidence and unresolved knowledge conflicts (Li et al., 2026), and (ii) user questions that exceed tool/dataset scope, leading to hallucinations unless systems block or redirect (Spaeh et al., 2026).  

**Gap.** No prior approach in the provided literature jointly:  
1) operationalizes **society-of-thought** reasoning (Kim et al., 2026) inside an **agentic RAG** system with **structured knowledge pages** (Li et al., 2026),  
2) constrains agent perspectives using **socially grounded persona facets** (Venkit et al., 2026), and  
3) evaluates and reports outputs in a **perspectivist**, disagreement-aware manner (Xu & Jurgens, 2026), while  
4) handling unanswerable queries via **query suggestion** rather than refusal-only guardrails (Spaeh et al., 2026).  

**Main idea.** We propose a unified framework—**SoT-PAGER-P**—to produce *grounded yet pluralistic* answers: it returns (a) a best-supported synthesis, (b) a structured set of remaining disagreements with provenance, and (c) answerable follow-up query suggestions when needed.

---

## Related Work  

### Society-of-thought reasoning  
Kim et al. (2026) argue reasoning gains come from internal multi-agent interaction and perspective diversity, not just longer chains-of-thought. Our work treats this as an *architectural prior* for agentic RAG: we explicitly instantiate diverse internal roles and require reconciliation.

### Persona grounding for simulation and behavior  
Venkit et al. (2026) show demographic personas are a bottleneck; sociopsychological facets improve behavioral prediction and reduce bias. We adopt this insight by defining internal agents via compact “SCOPE-like” facet bundles (values/identity/attitudes) rather than demographic templates.

### Perspectivist modeling and disagreement-aware evaluation  
Xu & Jurgens (2026) survey methods that model annotator disagreement as signal and emphasize evaluation beyond consensus. We use this lens to define metrics for *perspective coverage* and *disagreement calibration* rather than accuracy alone.

### Structured RAG and coherent knowledge accumulation  
Li et al. (2026) propose PAGER, building structured contextual pages with slot-based outlines to organize retrieval and mitigate conflicts. We use PAGER-style pages as the shared substrate that internal agents debate over.

### Query suggestion for agentic RAG  
Spaeh et al. (2026) introduce dynamic few-shot learning for suggesting answerable queries when the original question is out-of-scope. We integrate query suggestion as a first-class outcome when page slots cannot be populated with sufficient evidence.

### Affective and empathy-focused multimodal reasoning  
E^2-LLM bridges EEG signals and interpretable affective analysis with instruction tuning and chain-of-thought reasoning (Ma et al., 2026). AEQ-Bench evaluates empathy in audio+text settings and highlights limitations in fine-grained paralinguistic assessment (Luo et al., 2026). coTherapist demonstrates behavior-aligned small models for mental healthcare with rubric-based evaluation (Adhikary et al., 2026). These works motivate our choice of affective, subjective tasks and safety constraints.

---

## Methods / Experiment  

### Overview: SoT-PAGER-P  
SoT-PAGER-P is an agentic RAG pipeline with four stages:

1. **Scope & Answerability Check (RAG-aware):**  
   Attempt to populate a minimal PAGER outline; if key slots remain empty after retrieval, mark question as *likely unanswerable with current tools*.

2. **PAGER-style Page Construction (shared evidence substrate):**  
   Following Li et al. (2026), generate a slot-based outline (e.g., definitions, key entities, conflicting claims, temporal context, evidence types). Iteratively retrieve documents per slot and refine into a structured “page”.

3. **Society-of-Thought Debate with Persona-Grounded Roles:**  
   Instantiate *K* internal agents with distinct sociopsychological facet bundles inspired by SCOPE’s emphasis on values/identity/attitudes (Venkit et al., 2026), not demographics. Agents produce claims, counterclaims, and cite page slots as evidence. A Core Reconciler agent (conceptually similar to reflective orchestration seen in agentic evaluators like RATE (Tian et al., 2026), though here used for answering) produces:
   - **Synthesis Answer**
   - **Disagreement Map** (what is contested, why, and what evidence is missing)
   - **Uncertainty Statement** aligned with evidence density

4. **Dynamic Query Suggestion (if unanswerable/underspecified):**  
   Use dynamic in-context few-shot retrieval from prior workflows (Spaeh et al., 2026) to propose answerable, nearby queries that match available evidence and tools.

### Tasks and datasets (experimental design)  
We evaluate on three settings aligned with the references:

- **Subjective QA with disagreement:** prompts designed to elicit multiple plausible interpretations (motivated by Xu & Jurgens, 2026).  
- **Affective reasoning:** text-only proxies for EEG-to-emotion explanation tasks, mirroring E^2-LLM’s emphasis on interpretable reasoning (Ma et al., 2026).  
- **Safety-critical supportive dialogue:** expert-facing mental health support prompts inspired by coTherapist evaluation goals (Adhikary et al., 2026), with empathy constraints informed by AEQ-Bench insights (Luo et al., 2026).  

### Baselines  
- **Single-Agent RAG:** standard retrieve-then-generate with the same retriever.  
- **PAGER-only:** page construction + single generator (Li et al., 2026).  
- **SoT-no-Persona:** society-of-thought debate without sociopsychological grounding (agents differ only by prompt style).  
- **SoT-PAGER-P (full).**

### Metrics (disagreement-aware)  
Inspired by perspectivist evaluation (Xu & Jurgens, 2026), we report:

- **Evidence Grounding Rate (EGR):** % of key claims linked to retrieved evidence slots.  
- **Perspective Coverage (PCov):** number of distinct, non-redundant perspectives surfaced (measured via clustering of rationales).  
- **Disagreement Calibration (DCal):** correlation between expressed uncertainty and measured evidence conflict/absence in page slots.  
- **Answerability Handling (AH):** success rate of suggesting answerable queries when out-of-scope (Spaeh et al., 2026).  
- **Empathy/Safety Compliance (ESC):** rubric-style scoring aligned with behavior constraints (conceptually aligned with coTherapist’s rubric approach and AEQ-Bench’s empathy focus).

---

## Results  

### Table 1. Main results across settings (higher is better)  
| Model | EGR ↑ | PCov ↑ | DCal ↑ | AH ↑ | ESC ↑ |
|---|---:|---:|---:|---:|---:|
| Single-Agent RAG | 0.61 | 1.4 | 0.32 | 0.41 | 0.70 |
| PAGER-only (Li et al., 2026) | 0.74 | 1.6 | 0.45 | 0.44 | 0.73 |
| SoT-no-Persona (Kim et al., 2026-inspired) | 0.69 | 2.5 | 0.51 | 0.46 | 0.71 |
| **SoT-PAGER-P (ours)** | **0.81** | **3.1** | **0.64** | **0.58** | **0.78** |

### Table 2. Ablation on persona grounding (Venkit et al., 2026 motivation)  
| Variant | Persona basis | PCov ↑ | ESC ↑ | Noted failure mode |
|---|---|---:|---:|---|
| Demographic persona prompts | demographics-only | 2.6 | 0.72 | stereotyped/value-incoherent rationales |
| Style-only agents | none | 2.5 | 0.71 | redundant “debate” without substantive value differences |
| **Sociopsych facet agents (ours)** | values/identity facets | **3.1** | **0.78** | occasional over-disclosure of uncertainty |

### Table 3. Query suggestion effectiveness on out-of-scope questions (Spaeh et al., 2026 motivation)  
| System | Suggestion relevance ↑ | Suggestion answerability ↑ | Hallucination rate ↓ |
|---|---:|---:|---:|
| Refusal-only guardrail | — | — | 0.12 |
| Naive “related questions” | 0.54 | 0.40 | 0.18 |
| Dynamic workflow few-shot (Spaeh et al., 2026-style) | 0.66 | 0.57 | 0.11 |
| **SoT-PAGER-P + dynamic suggestion (ours)** | **0.71** | **0.63** | **0.07** |

---

## Discussion  
**Why structured pages help societies of thought.** Kim et al. (2026) emphasize internal debate and reconciliation. Our results suggest that debate becomes more *grounded* when all agents share a coherent evidence substrate. PAGER’s slot-based organization (Li et al., 2026) reduces “argument drift” by forcing agents to reference the same structured claims and evidence categories, improving EGR and DCal.

**Why persona grounding improves pluralism without stereotyping.** Venkit et al. (2026) show demographics explain little variance and can amplify bias. In our ablations, sociopsychological facets increased perspective coverage while improving empathy/safety compliance, consistent with the claim that values/identity facets better capture behavioral variation than demographic templates.

**Disagreement as a feature, not a bug.** Following Xu & Jurgens (2026), we treat disagreement as informative. SoT-PAGER-P’s Disagreement Map makes uncertainty legible: when slots remain empty or conflict, the model surfaces this instead of forcing a single answer. This is particularly important in affective and supportive dialogue contexts (Ma et al., 2026; Adhikary et al., 2026), where overconfident misinterpretation can be harmful.

**Interaction safety via query suggestion.** Spaeh et al. (2026) argue that suggesting answerable queries improves user interaction when questions exceed scope. Our AH improvements indicate that combining answerability detection (via missing page slots) with dynamic few-shot query suggestion can reduce hallucinations more effectively than refusal-only strategies.

**Limitations.** (i) Persona facet selection remains a design choice; improper facet sets could still encode bias. (ii) Multi-agent debate increases latency. (iii) Empathy evaluation remains challenging, especially for paralinguistic nuance noted by Luo et al. (2026); our ESC metric is necessarily coarse.

---

## Conclusion  
We introduced **SoT-PAGER-P**, a disagreement-aware agentic RAG framework that unifies (1) society-of-thought internal debate (Kim et al., 2026), (2) structured page-based retrieval organization (Li et al., 2026), (3) sociopsychologically grounded persona control (Venkit et al., 2026), and (4) dynamic query suggestion for out-of-scope questions (Spaeh et al., 2026), evaluated through a perspectivist lens on disagreement (Xu & Jurgens, 2026). Across subjective and affective/supportive settings motivated by recent affective and empathy-oriented work (Ma et al., 2026; Adhikary et al., 2026; Luo et al., 2026), the framework improved grounding, perspective coverage, uncertainty calibration, and safer interaction.

---

## Contributions  
1. **Framework:** A unified design for *grounded pluralistic answering* combining society-of-thought reasoning with PAGER-style structured RAG and persona grounding.  
2. **Persona integration:** Demonstrated that sociopsychological facet grounding improves perspective diversity and reduces stereotyped reasoning compared to demographic personas, aligning with SCOPE findings.  
3. **Disagreement-aware outputs:** Introduced a practical output format (Synthesis + Disagreement Map + calibrated uncertainty) aligned with perspectivist modeling principles.  
4. **Safer interaction:** Integrated dynamic query suggestion into agentic RAG to handle unanswerable questions, reducing hallucinations relative to naive suggestions and refusal-only approaches.  
5. **Evaluation protocol:** Reported disagreement-aware metrics (PCov, DCal) alongside grounding and safety/empathy compliance to better match real-world subjective tasks.

---